% Options for packages loaded elsewhere
\PassOptionsToPackage{unicode}{hyperref}
\PassOptionsToPackage{hyphens}{url}
\documentclass[
]{article}
\usepackage{xcolor}
\usepackage{amsmath,amssymb}
\setcounter{secnumdepth}{-\maxdimen} % remove section numbering
\usepackage{iftex}
\ifPDFTeX
  \usepackage[T1]{fontenc}
  \usepackage[utf8]{inputenc}
  \usepackage{textcomp} % provide euro and other symbols
\else % if luatex or xetex
  \usepackage{unicode-math} % this also loads fontspec
  \defaultfontfeatures{Scale=MatchLowercase}
  \defaultfontfeatures[\rmfamily]{Ligatures=TeX,Scale=1}
  % unicode-math maps \setminus to U+29F5 (REVERSE SOLIDUS OPERATOR), which
  % latinmodern-math.otf lacks (tectonic/lualatex emit "Missing character
  % U+29F5"). Remap to U+2216 (SET MINUS), matching the main paper (commit
  % 9f0c391) and pdfTeX/amssymb rendering. Output-preserving.
  \AtBeginDocument{\renewcommand{\setminus}{\mathbin{\char"2216}}}
\fi
\usepackage{lmodern}
\ifPDFTeX\else
  % xetex/luatex font selection
\fi
% Use upquote if available, for straight quotes in verbatim environments
\IfFileExists{upquote.sty}{\usepackage{upquote}}{}
\IfFileExists{microtype.sty}{% use microtype if available
  \usepackage[]{microtype}
  \UseMicrotypeSet[protrusion]{basicmath} % disable protrusion for tt fonts
}{}
\makeatletter
\@ifundefined{KOMAClassName}{% if non-KOMA class
  \IfFileExists{parskip.sty}{%
    \usepackage{parskip}
  }{% else
    \setlength{\parindent}{0pt}
    \setlength{\parskip}{6pt plus 2pt minus 1pt}}
}{% if KOMA class
  \KOMAoptions{parskip=half}}
\makeatother
\setlength{\emergencystretch}{3em} % prevent overfull lines
\providecommand{\tightlist}{%
  \setlength{\itemsep}{0pt}\setlength{\parskip}{0pt}}
\usepackage{bookmark}
\IfFileExists{xurl.sty}{\usepackage{xurl}}{} % add URL line breaks if available
\urlstyle{same}
\hypersetup{
  hidelinks,
  pdfcreator={LaTeX via pandoc}}

\author{}
\date{}

\begin{document}

\textbf{RNR Coding --- Summary of Formal Results}

\emph{Companion to main paper, v0.8}

Pavlo Ratushnyi --- May 2026

\textbf{Purpose}

This document collects all formal results of the RNR Coding framework in
one place, with a one-line statement, a one-line proof sketch, and a
one-line interpretation per result. Use it for talk preparation, defense
rehearsal, or as a quick-reference card when answering reviewer
questions. Section numbers refer to the main paper.

\textbf{Notation reminder}

\begin{quote}
$X$: source block. $Y = M(C)$: decoder-reproducible side information
from public context $C$. $Q$: predicted probability field. $U$:
predicted uncertainty/repair-cost field. $R$: repair stream.
$L(\cdot)$: bit length. $H(\cdot)$: Shannon entropy. $\mathrm{TC}$:
total correlation (multi-information). $\varepsilon_n$:
entropy-coder redundancy.
\end{quote}

\textbf{Section 4 --- Type-I: Stream Repair Coding}

\textbf{Theorem 4.1 (Type-I achievability) --- {[}achievability{]}}

\begin{quote}
\emph{Statement.} $\mathbb{E}[L_I(X)] \leq L(M_I) + N \cdot H(P_E) +
N \cdot \ell_{\text{cond}} + \varepsilon_N + L(V) + L(\Theta)$, where
$P_E$ is the empirical Code12 error-mask class distribution.

\emph{Proof.} Direct composition of class entropy (Shannon-coding
bound on the mask classes), within-class payload cost, and constant
overheads.

\emph{Why it matters.} Establishes that Type-I attains conditional
source-coding rate plus $O(1)$ overhead, modulo the entropy-coder
redundancy. Cite when claiming Type-I is asymptotically competitive.
\end{quote}

\textbf{Theorem 4.2 (Type-I converse) --- {[}lower bound{]}}

\begin{quote}
\emph{Statement.} $\mathbb{E}[L(C(X) \mid Y)] \geq H(X \mid Y)$ for
any uniquely decodable code with deterministic decoder reproducing
$Y$.

\emph{Proof.} Shannon's noiseless source coding applied to the
conditional distribution $P(X \mid Y)$.

\emph{Why it matters.} Pairs with Theorem 4.1: the achievability
bound is order-optimal. No coder using only $Y$ can do better than
$H(X \mid Y)$. When pressed on ``could you do better than this
rate?'' the answer is: no, by this converse.
\end{quote}

\textbf{Remark 4.3a (Fixed-size Code12 is constant-size) --- {[}disclaimer{]}}

\begin{quote}
\emph{Statement.} For the deployed Code12 parameters $n_0 = 256$,
$m_0 = 12$, optimal $E_{12}$ is a finite enumeration over
$\binom{2^{12}}{256} \cdot 256!$ injective maps --- a constant
independent of $|F|$. Hence the practical problem is polynomial in
the formal asymptotic sense; NP-hardness applies only when $n$ or
$m$ is part of the input.

\emph{Why it matters.} Honest disclaimer. The earlier draft claimed
NP-hardness for fixed $(256, 12)$; that claim is formally vacuous and
has been removed.
\end{quote}

\textbf{Theorem 4.3 (Varying-parameter NP-hardness) --- {[}complexity{]}}

\begin{quote}
\emph{Statement.} For varying $(n, m)$ with $n \leq 2^m$, symmetric
row-stochastic $F \in \mathbb{Q}_{\geq 0}^{n \times n}$, and rational
$B$, deciding whether there is an injective $E : [n] \to \{0, 1\}^m$
with $\sum_{i,j} F_{ij} d_H(E(i), E(j)) \leq B$ is NP-hard.

\emph{Proof.} Reduce from Wagner-Corneil (1990) tree-to-hypercube
subgraph embedding. Given a tree $T$, build symmetric row-stochastic
$F$ with $F_{uv} = 1/\Delta$ on tree edges, $F_{uu} = 1 -
\deg_T(u)/\Delta$, and threshold $B = 2(n-1)/\Delta$. By injectivity
the objective is $\geq B$ with equality iff every tree edge has
Hamming distance one, iff $T$ embeds as a subgraph of $Q_m$.

\emph{Why it matters.} Replaces the broken QAP-Hamming-embedding
proof with a sound reduction from a Hamming-native NP-hard problem.
Justifies heuristics for $E_{12}$ in the asymptotic regime.
\end{quote}

\textbf{Remark 4.3b (Asymmetric confusion matrices) --- {[}reduction{]}}

\begin{quote}
\emph{Statement.} For any nonnegative $F$, $\sum F_{ij} d_H(E(i),
E(j)) = \sum \frac{F_{ij}+F_{ji}}{2} d_H(E(i), E(j))$ because $d_H$
is symmetric. Thus an arbitrary asymmetric empirical confusion
matrix induces the same labeling objective as its symmetric part.

\emph{Why it matters.} The symmetric-input restriction in Theorem 4.3
is not a real restriction.
\end{quote}

\textbf{(Theorem 4.4 deleted.)} The earlier APX-hardness claim
inherited the broken QAP-Hamming-embedding reduction. No clean
constant-gap reduction from an APX-hard source into the
Hamming-realizable geometry is currently known; constant-gap
inapproximability for $E_{12}$ is listed as an open problem.

\textbf{Section 5 --- Type-II: Positional Symbol Coding}

\textbf{Lemma 5.1 (factorized equivalence) --- {[}equivalence{]}}

\begin{quote}
\emph{Statement.} Under the factorized regime, per-position
arithmetic coding under $Q$ and per-position Type-I arithmetic
coding under the equivalent sequential conditional $Q$ are
bit-equivalent.

\emph{Proof.} Both coders are arithmetic coders driven by the same
probability vector at each position; redundancy is the constant
$\varepsilon_N$.

\emph{Why it matters.} Shows that factorized Type-II is not
operationally distinct from Type-I. The interesting Type-II case is
non-factorized, which Theorem 5.4 then exhibits.
\end{quote}

\textbf{Theorem 5.2 (multinomial baseline) --- {[}combinatorial
bound{]}}

\begin{quote}
\emph{Statement.} The number of valid 16-way partitions of
$\{1, \ldots, N\}$ with given counts is the multinomial
$M(N; k_0, \ldots, k_{15})$; enumerative coding describes a partition
in $\log_2 M$ bits.

\emph{Proof.} Combinatorial counting: partitions correspond
bijectively to multiset labelings.

\emph{Why it matters.} The correct baseline for positional coding.
Earlier drafts used a looser sum-of-binomials baseline that
overstated gains; this is the v0.3 correction.
\end{quote}

\textbf{Theorem 5.3 (support-set advantage with escape) ---
{[}achievability{]}}

\begin{quote}
\emph{Statement.} If the predictor supplies candidate supports
$S_v$ with $\alpha_v$ escapes, the partition cost is at most
$\sum_v [\log_2 \binom{s_v}{k_v - \alpha_v} + \log_2 \binom{N - s_v}{\alpha_v}] + L(\text{counts}) + \varepsilon_N$.

\emph{Proof.} Decompose $M_v$ as the disjoint union $M_v \cap S_v$
plus $M_v \setminus S_v$; encode each by enumerative rank.

\emph{Why it matters.} Quantifies the gain from predicted support
information with explicit escape cost. Lets a reviewer see exactly
what the predictor must achieve for Type-II to win.
\end{quote}

\textbf{Theorem 5.4 (semi-non-factorized separation, linear-code form)
--- {[}separation{]}}

\begin{quote}
\emph{Statement.} There exists a distribution $D_N$ over
$A_{16}^N$ (uniform marginals, image of a $(4-c)N$-to-$4N$ linear
code) such that any factorized arithmetic coder pays $\geq 4N$ bits
while a semi-non-factorized Type-II coder achieves
$(4-c)N + O(\log N)$ bits. Separation is $\Theta(N)$.

\emph{Proof.} Linear-code construction over $\mathrm{GF}(2)$: a
generator matrix $G : \mathrm{GF}(2)^{(4-c)N} \to \mathrm{GF}(2)^{4N}$
whose four columns per nibble are linearly independent yields exactly
uniform per-nibble marginals; membership in the codeword set fits in
$(4-c)N$ bits.

\emph{Why it matters.} Settles whether semi-non-factorized Type-II
genuinely beats Type-I: yes, by $\Theta(N)$ bits in the worst case.
Strongest separation result in the paper.
\end{quote}

\textbf{Theorem 5.5 (multi-information separation, general source) ---
{[}separation, general{]}}

\begin{quote}
\emph{Statement.} For any source $D_N$, factorized-vs-joint
separation equals the total correlation
$\mathrm{TC}(D_N) = \sum_i H_{\text{marg}}(X_i) - H(D_N)$.

\emph{Proof.} Both rates are Shannon-optimal for their respective
coders; their difference is $\mathrm{TC}$ by definition.

\emph{Why it matters.} Generalizes 5.4 from the linear-code
construction to arbitrary sources. The right unifying object for
understanding when Type-II wins is total correlation.
\end{quote}

\textbf{Remark 5.6 (manifold-hypothesis quantification of TC) ---
{[}conditional quantitative claim{]}}

\begin{quote}
\emph{Statement.} Under the manifold hypothesis (natural data has
intrinsic dimension $d \ll N$),
$\mathrm{TC}(D_N) = \Theta(N) - O(d \log N)$, supported by published
intrinsic-dimension estimates for images, text, and structured
records.

\emph{Proof.} Combination of Theorem 5.5 with empirical
intrinsic-dimension results from Carlsson et al., Levina-Bickel,
Tenenbaum et al.

\emph{Why it matters.} Bridges the conditional theorem to natural data.
Use when a reviewer asks why the framework is relevant beyond contrived
families.
\end{quote}

\textbf{Remark 5.7 (Support optimization is open) --- {[}status{]}}

\begin{quote}
\emph{Statement.} The natural marginal-$Q$-prefix heuristic is not
provably optimal for fixed-size Type-II support selection: even when
$A_s(\alpha) = \log_2 \binom{s}{k - \alpha} + \log_2 \binom{N - s}{\alpha}$
is non-decreasing on feasible escape counts, the expected cost depends
on the full joint law of $M_v$, not only marginals $Q_v(i)$.

\emph{Counterexample.} Concrete falsification at $N = 7$, $k = 2$,
$s = 2$ in the verification script
\texttt{scripts/verify/prop\_5\_7\_counterexample.py}: top-$Q$ support
$\{2,4\}$ costs 2.658 bits; support $\{0,4\}$ costs 2.126 bits.

\emph{Why it matters.} Replaces a previously-asserted prefix-optimality
proposition. The arithmetic-coding benchmark (cross-entropy under $Q$,
equal to $H(X \mid Q)$ only when $Q$ is the true conditional law)
remains a separate result.
\end{quote}

\textbf{Section 6 --- Type-III: Tensor Entropy-Field Coding}

\textbf{Theorem 6.1 (approximate-oracle mode selection) --- {[}upper
bound{]}}

\begin{quote}
\emph{Statement.} If $|U_m(t) - L_m(t)| \leq \varepsilon_t$ for all
$m$, then $L_{\hat{m}(t)}(t) \leq L_{m^\star(t)}(t) + 2\varepsilon_t$,
where $\hat{m}$ is the predicted argmin and $m^\star$ the oracle
argmin.

\emph{Proof.} Three triangle-inequality steps over the prediction
error bound and the definition of argmin.

\emph{Why it matters.} Defends Type-III-A composition: a calibrated
$U$ lets the framework approach oracle mode selection. The
$2\varepsilon$ factor is the key constant.
\end{quote}

\textbf{Theorem 6.2 ($2\varepsilon$ bound is tight) ---
{[}matching lower bound{]}}

\begin{quote}
\emph{Statement.} For any $\varepsilon > 0$ and $\delta > 0$, an
adversarial two-mode configuration achieves
$L_{\hat{m}} \geq L_{m^\star} + 2\varepsilon - \delta$.

\emph{Proof.} Construct $L_1 = \delta/2$,
$L_2 = 2\varepsilon - \delta/2$, $U_1 = \delta/2 + \varepsilon$,
$U_2 = \varepsilon - \delta/2$; the argmin selects the inferior
mode by exactly $2\varepsilon - \delta$.

\emph{Why it matters.} Pairs with 6.1: the $2\varepsilon$ factor
cannot be improved without additional assumptions on the predictor's
error structure.
\end{quote}

\textbf{Theorem 6.2b (Type-III-A absolute-rate converse) ---
{[}converse{]}}

\begin{quote}
\emph{Statement.} Any uniquely-decodable Type-III-A code obeys
$\mathbb{E}[L_{\mathrm{IIIA}}] \geq H(X \mid Y)$. Under the per-tile-UD
coding hypothesis, the causal-adaptive floor
$R_{\mathrm{tile}} = \sum_t \mathbb{E}_{C_t}[\min_m
\mathbb{E}[L_m(t) \mid C_t]]$ satisfies $R_{\mathrm{tile}} =
H(X \mid Y)$ iff, per tile and per context, the selected mode is
expressive AND its context $C_t$ is causally sufficient; otherwise the
slack is the context deficit, at most the inter-tile total correlation
$\mathrm{TC}(X_1,\ldots,X_T \mid Y)$ (the block-independent worst case),
and zero under causal-context modes.

\emph{Proof.} McMillan plus the conditional Shannon bound for (6.2b.1);
for the floor, $\min_m \mathbb{E}[L_m(t) \mid C_t] \geq H(X_t \mid C_t)
\geq H(X_t \mid X_{<t}, Y)$, summed via the chain rule
$\sum_t H(X_t \mid X_{<t}, Y) = H(X \mid Y)$.

\emph{Why it matters.} Completes the achievability--converse matrix:
with 6.2 (relative) and 6.2b (absolute), Type-III-A joins Types I, II,
III-B, III-C in having both a coding bound and a matching converse. The
operative cost is the \emph{causal-adaptive} floor --- a context-free
per-tile commit over-estimates it, the hindsight envelope under-estimates
it (and is non-UD).
\end{quote}

\textbf{Theorem 6.3 (byte-tensor mutual information advantage) ---
{[}achievability{]}}

\begin{quote}
\emph{Statement.} Byte-tensor coding's advantage over independent
nibble coding equals $I(H ; L \mid C)$, the conditional mutual
information between high and low nibbles given context.

\emph{Proof.} Chain rule of entropy:
$H(H, L \mid C) = H(H \mid C) + H(L \mid H, C)$; difference from
independent coding is exactly $I(H ; L \mid C)$.

\emph{Why it matters.} Quantifies exactly when Type-III-B wins. Use
when a reviewer asks: ``what does joint byte modeling actually buy
you?''
\end{quote}

\textbf{Theorem 6.4 (hierarchical decomposition is invariant) ---
{[}invariance{]}}

\begin{quote}
\emph{Statement.} Any rooted binary tree decomposition of the
256-byte alphabet, used as the order of conditional encoding,
achieves expected length $H(B \mid C)$. The tree shape does not
affect the rate.

\emph{Proof.} Chain rule of entropy is path-independent: reordering
conditional entropies along any spanning tree gives the same sum
$H(B \mid C)$.

\emph{Why it matters.} The choice between high-low, low-high, or
arbitrary tree decompositions is a calibration question, not a rate
question. Information-theoretically free.
\end{quote}

\textbf{Theorem 6.5 (MDL token admission) --- {[}optimization
criterion{]}}

\begin{quote}
\emph{Statement.} Token $\tau$ is admitted to the dictionary iff
$f \cdot [8\ell - \mathrm{code\_cost}(\tau)] > \mathrm{dictionary\_cost}(\tau) + f \cdot \mathrm{parse\_overhead}(\tau)$.

\emph{Proof.} Net code-length change from adding $\tau$ is the LHS
minus RHS; $\tau$ saves bits iff this is positive.

\emph{Why it matters.} Standard MDL admission rule, placed in the
framework to make explicit that not every byte pattern earns
dictionary status.
\end{quote}

\textbf{Theorem 6.6 (bits-back integration) --- {[}equivalence with KL
accounting{]}}

\begin{quote}
\emph{Statement.} Under bits-back ANS, the Type-III-C repair length
$L(R \mid Y) = L_{\text{prior}}(z \mid Y) + L_{\text{cond}}(X \mid z, Y) - L_{\text{aux}}(z)$
satisfies
$\mathbb{E}[L(R \mid Y)] = H(X \mid Y) + \mathrm{KL}(q \,\Vert\, P_{\text{posterior}})$.

\emph{Proof.} Three-stream accounting: encode latent $z$ under
prior, data under conditional, recover auxiliary bits used to sample
$z$ under $q$; expectation reduces via joint factorization to
marginal $H(X \mid Y)$ plus posterior KL.

\emph{Why it matters.} Fits bits-back coding into the framework with
no overclaim. When $q$ matches the posterior, rate equals the
Slepian-Wolf bound. Cites BB-ANS and Bit-Swap.
\end{quote}

\textbf{Theorem 6.7 (joint dictionary + predictor NP-hard) ---
{[}complexity{]}}

\begin{quote}
\emph{Statement.} Joint dictionary-and-predictor learning is
NP-hard, even when $M$ is fixed to the trivial constant predictor and
$D$ is restricted to the recursive (straight-line grammar) dictionary
class.

\emph{Proof.} Polynomial-time reduction from the Smallest Grammar
Problem (Charikar et al. 2005): under the constant predictor on the
recursive-$D$ class, the Type-III-C objective reduces to grammar
size; the bijection between dictionaries and straight-line grammars
preserves optimality.

\emph{Why it matters.} Defends the heuristic nature of practical
dictionary learning (BPE, MDL-greedy). Optimal learning is not
achievable in polynomial time.
\end{quote}

\textbf{Theorem 6.8 (joint dictionary inapproximable) --- {[}strong
complexity{]}}

\begin{quote}
\emph{Statement.} Joint dictionary-and-predictor learning is
APX-hard: some constant $\gamma > 1$ exists such that no polynomial-
time $\alpha$-approximation exists for $\alpha < \gamma$, unless
$P = NP$ (the explicit numerical $\gamma$ is left open).

\emph{Proof.} The SGP reduction in 6.7 transfers Charikar et al.'s
APX-hardness of SGP up to the encoding-convention slack between
symbol-count and bit-count objectives.

\emph{Why it matters.} Strengthens 6.7: even constant-factor
approximation is hard. Pairs with 4.4 to give the framework a complete
complexity classification of its two NP-hard problems.
\end{quote}

\textbf{Lemma 6.9 + Remarks 6.9a/c/d + Corollary 6.9b (continuous
Type-III-C support selection) --- {[}complexity dichotomy{]}}

\begin{quote}
\emph{Statement.} For the Gaussian-field (bits-back limit) Type-III-C
source, the per-block residual rate is the Schur-complement log-det
$\tfrac12\log_2\det\Sigma_{S\mid\bar S}$. With the field \emph{fixed},
minimising it over size-$s$ supports is maximum-entropy sampling,
APX-hard with constant $5/4$ under $\mathrm{P}\neq\mathrm{NP}$
(Lemma 6.9, via Ko--Lee--Queyranne 1995 / Ohsaka 2022;
\textsc{Partial}). With the predictor \emph{free and amortised} the
same objective is in $\mathrm{P}$ via a sample-covariance MLE plug-in
(Remark 6.9a). The crossover budget is delimited in Remark 6.9c, the
$\Sigma\succeq I$ calibration shown necessary in Remark 6.9d, and an
operational additive-regret encoder synthesised in Corollary 6.9b.

\emph{Proof.} Bits-back rate (Thm 6.6) $=$ conditional differential
entropy $=$ Schur-complement log-det (Cover--Thomas Thm 8.4.1); MESP
inapproximability for the fixed field; no-free-lunch identity
$h(X_S)=h(X_S\mid Y)+I(X_S;Y)$ for the amortised collapse.

\emph{Why it matters.} Locates the genuine \emph{non-count} hardness
gap that provably cannot live in Type-II (it does not dilute under the
affine identity) in the continuous Type-III-C objective, while honestly
bounding it to the fixed-field/single-block regime --- the continuous
mirror of OP2(b). Cite when asked whether any RNR support-selection
problem is genuinely APX-hard for a real cost.

\emph{Discrete transfer (Remark 7.15a).} The same framework ports to the
discrete autoregressive neural-LM field of Theorem 7.15: discrete entropy
$H(X_T)$ is submodular with \emph{no} calibration (Fujishige 1978;
Remark 6.9d's $\Sigma\succeq I$ floor vanishes), greedy keeps the
$(1-1/e)$ additive-regret guarantee, and a value-oracle complexity split
appears with \emph{no} Gaussian analogue --- order-1-arbitrary and
bounded-order-contiguous poly, but the \emph{exact} oracle is \#P-hard for
general fields (Roth 1996). So the deployed discrete setting carries two
obstacles (NP-hard selection \emph{and} \#P-hard value oracle) where the
Gaussian theory carries one.
\end{quote}

\textbf{Sections 7--10 --- Composition, hardness, verification,
reproducibility}

\textbf{Theorem 7.1 (composition correctness) --- {[}construction{]}}

\begin{quote}
\emph{Statement.} The Type-III-A composition with mode set $M$ and
deterministic $\hat{m} = \operatorname*{arg\,min}_m U_m(t)$ produces
a uniquely decodable codeword for $X$.

\emph{Proof.} Each step in the decoding pipeline (metadata read, $U$
reconstruction, $\hat{m}$ recovery, per-tile repair decoding) is
uniquely determined.

\emph{Why it matters.} Establishes that the framework's hierarchical
composition is well-formed. No ambiguity in the decoding pipeline.
\end{quote}

\textbf{Theorems 7.2--7.28 --- RNR with byte-granular random access}

The §7.2--7.28 expansion delivers a 26-theorem framework (plus the
7.6$'$/7.27$'$ refinements and the 7.11/7.27c/7.27e/7.27f corollaries)
bridging the abstract RNR formalism (§§3--6) to a deployable
neural-lossless-compression pipeline. Headline results, organized by
theme:

\textbf{Theorem 7.2 (byte-granular RA architecture) --- {[}construction{]}}

\begin{quote}
\emph{Statement.} Sub-block-based RNR with sync points every $K$
bytes admits byte-granular random-access decoding at cost
$O(\log L(R) + K \cdot \mathrm{cost}(M'))$ per query and total
encoded length $N \cdot H_{M'}(X) + m \cdot \delta_\infty + O(m \cdot p)$
under conventions C1--C3 (stationary, $W_N$-bounded predictor,
$\eta$-floor AC at $p$-bit precision, $\delta_\infty$ per-position
Crutchfield--Feldman excess entropy).

\emph{Why it matters.} The architectural foundation for everything in
§§7.3--7.24. Establishes the storage/access trade-off via $K$, with
$K = \Theta(\sqrt{N})$ optimal balance (Theorem 7.5).
\end{quote}

\textbf{Theorem 7.5 (optimal $K^\ast = \Theta(\sqrt{N})$) ---
{[}optimization{]}}

\begin{quote}
\emph{Statement.} For balanced storage/access workloads with per-query
random access cost $\alpha \cdot K \cdot \mathrm{cost}(M')$ and total
storage cost $\beta \cdot m \cdot \delta_\infty + N H_{M'}$ for
$m = N/K$, the cost-minimizing sub-block size is
$K^\ast = \sqrt{\beta \delta_\infty / (\alpha \cdot Q)}$ giving
total $\Theta(\sqrt{N})$ in $m$ and per-query.

\emph{Why it matters.} Determines the sweet-spot sub-block size for
production: too small $K$ blows up storage overhead; too large $K$
blows up per-query latency. The $\sqrt{N}$ balance is universal under
Theorem 7.2's hypotheses.
\end{quote}

\textbf{Theorems 7.6--7.10 --- Source-coding extensions
{[}composition{]}}

\begin{quote}
T7.6: universal online RNR via Krichevsky--Trofimov estimator with
regret $O(\sqrt{N \log |\Sigma|})$ under unknown source. T7.7:
successive refinement RNR via Equitz--Cover 1991 (multi-resolution
layering). T7.8: functional compression via Orlitsky--Roche 2001
characteristic-graph entropy. T7.9: lossy RNR with distortion budget
$D$ via Shannon $R(D)$ (corrected from earlier ``Westerlund-Wyner''
fabrication to Steinberg--Merhav 2004). T7.10: Slepian--Wolf /
Wyner--Ziv side-information composition.

\emph{Why it matters.} Demonstrates RNR's compositional reach beyond
the basic random-access architecture. Covers the classical 5
source-coding extensions in a unified framework.
\end{quote}

\textbf{Theorem 7.6$'$ (tight $d/2$ minimax regret) --- {[}sharp
constant{]}}

\begin{quote}
\emph{Statement.} For a $d$-dimensional smooth parametric source
class, the universal online RNR predictor (Jeffreys-prior mixture /
per-context KT) attains regret $\tfrac{d}{2}\log_2 N + C + o(\log N)$,
and this constant is tight: for a.e.\ source parameter no online RNR
scheme beats $\tfrac{d}{2}(1-\varepsilon)\log_2 N$.

\emph{Proof.} Upper bound: Clarke--Barron / Krichevsky--Trofimov
asymptotic redundancy. Lower bound: Rissanen's $(d/2)\log N$ minimax
converse (Shtarkov pointwise-redundancy form).

\emph{Why it matters.} Sharpens T7.6's $d\log N$ to its true leading
constant $d/2$ --- exactly half --- and certifies it unimprovable.
The honest answer to ``what is the no-pre-training overhead?''.
\end{quote}

\textbf{Corollary 7.11 (lossy + ECC = lossless) --- {[}observation{]}}

\begin{quote}
\emph{Statement.} For any lossy RNR at distortion budget $D$ with
predictor $M'$ and per-symbol error probability $p$ under $D$, the
composition with an outer error-correcting code at rate
$R_{\mathrm{ecc}} \geq h_2(p)$ recovers lossless RNR with rate
$R^{\mathrm{lossless}}_{\mathrm{RNR}} \leq R(D) + h_2(p)$.

\emph{Why it matters.} A clean observation that bridges T7.9 (lossy)
to baseline lossless RNR via a standard channel-coding sandwich.
The user-originated insight ``lossy + ecc always can be lossless''
formalized.
\end{quote}

\textbf{Theorem 7.14 (concentration of $L^{\mathrm{rep}}$) ---
{[}deviation bound{]}}

\begin{quote}
\emph{Statement.} Under hypothesis (5.8) exp-mixing with constant
$\rho$, the encoder output length $L^{\mathrm{rep}}_{\mathrm{RNR}}$
concentrates around its expectation with Doob-martingale tail
\[
\Pr\!\Bigl(|L^{\mathrm{rep}} - \mathbb{E}[L^{\mathrm{rep}}]| > t\Bigr)
\;\leq\;
2 \exp\!\Bigl(-\frac{t^2}{2 N (W_N + 1 + \rho/(1-\rho))^2 \log_2^2(1/\eta)}\Bigr).
\]

\emph{Proof.} Doob-martingale $M_n = \mathbb{E}[L^{\mathrm{rep}} \mid
\mathcal{F}_n] - \mathbb{E}[L^{\mathrm{rep}}]$ has differences bounded
by $(W_N+1)$-context plus $\rho/(1-\rho)$ source-mixing-tail
contributions; Azuma--Hoeffding.

\emph{Why it matters.} Strengthens the expectation bound (T7.2) to a
high-probability bound. The $(W_N+1+\rho/(1-\rho))$ factor was refined
in response to a codex counterexample (W=0 persistent Markov chain
defeats the naive $(W_N+1)$ bound under slow mixing).
\end{quote}

\textbf{Theorem 7.25 (variance-aware Bernstein concentration) ---
{[}sharper deviation bound{]}}

\begin{quote}
\emph{Statement.} Under exp-mixing, $L^{\mathrm{rep}}$ concentrates
with a Bernstein tail driven by the long-run variance
$v^2 \approx \sigma_\ell^2(2W_N+1)$ (per-position-cost std
$\sigma_\ell$) rather than the worst-case Lipschitz $B = \log_2(1/\eta)$
of T7.14 --- an approximate $500\times$ tightening on enwik9
(rigorous $\delta = 10^{-6}$ band $\approx 1.3\,\mathrm{GB} \to
\approx 2$--$4\,\mathrm{MB}$).

\emph{Proof.} Merlev\`ede--Peligrad--Rio 2009 Bernstein inequality for
$\alpha$-mixing sums; Davydov 1968 covariance inequality for the
long-run-variance proxy.

\emph{Why it matters.} The preferred concentration bound for
exp-mixing sources; closes most of the gap between the loose Azuma
worst case and the (non-rigorous) CLT estimate. Use for archive
pre-allocation sizing.
\end{quote}

\textbf{Theorem 7.28 (non-asymptotic Berry--Esseen CLT) ---
{[}two-sided distributional bound{]}}

\begin{quote}
\emph{Statement.} The CDF of the normalised repair length is within
$O(N^{-1/2}(\log N)^2)$ (Kolmogorov distance) of a Gaussian, giving
\emph{exact} two-sided confidence intervals
$L^{\mathrm{rep}} \in \mathbb{E}[L^{\mathrm{rep}}] \pm
z_{1-\delta/2}\sqrt{N v^2}$ valid up to the additive CDF error.

\emph{Proof.} Berry--Esseen for $\alpha$-mixing sequences (Sunklodas
type) with the $(\log N)^2$ factor from the $W_N$-window dependence.

\emph{Why it matters.} Completes the Azuma/Bernstein/Berry--Esseen
trilogy: T7.14 gives a rigorous worst-case tail, T7.25 a tight
variance-aware tail, T7.28 the full two-sided distributional
approximation that rigorously justifies the CLT confidence bands used
in the T7.15 enwik9 prediction.
\end{quote}

\textbf{Theorem 7.15 (end-to-end Chinchilla-on-enwik9 prediction) ---
{[}numerical synthesis{]}}

\begin{quote}
\emph{Statement.} Substituting Chinchilla-70B $H_{M'} \approx 0.664$
bits/byte (Delétang et al.\ 2024) into Theorems 7.2--7.5--7.14 for
$N = 10^9$ enwik9 bytes, $K = \sqrt{N} \approx 31{,}623$,
$\eta = 2^{-32}$, $\delta = 10^{-6}$ predicts:
\[
L^{\mathrm{rep}}_{\mathrm{RNR}} \approx 79\,\mathrm{MB}
\quad (\text{12.0$\times$ bitstream-rate compression}),
\]
with two concentration bands: rigorous Azuma worst-case
$\pm 1.4\,\mathrm{GB}$ (provably correct upper bound, hopelessly
loose for realistic transformers with $W_N = 2048$) and practical
CLT-based estimate $\pm 0.04\,\mathrm{MB}$ at $\delta = 10^{-6}$ (per
T7.14's ``Relation to CLT'' remark with empirical $\sigma^2_{\ell_1}
\approx 4$). The CLT estimate matches Delétang et al.\ Table 1
measurement exactly.

\emph{Caveat (side-information framing).} The 12$\times$ figure is the
\emph{bitstream rate} excluding the 140 GB Chinchilla-70B model itself.
For self-contained algorithmic-IT framing, archive expansion is
$\approx 140\times$ at single-document scale; amortization break-even
at $V^\ast \approx 152\,\mathrm{GB}$ of compressible text.

\emph{Why it matters.} The headline practical result: RNR can deliver
Chinchilla-class compression at byte-granular random access.
Verified by \texttt{scripts/verify/thm\_7\_15\_neural\_rnr\_end\_to\_end.py}.
\end{quote}

\textbf{Theorems 7.16--7.19 --- Adaptive, edit-tolerant, parallel,
streaming {[}engineering{]}}

\begin{quote}
T7.16: adaptive per-region $K_r$ for piecewise-stationary sources with
strict Cauchy--Schwarz improvement over fixed $K$. T7.17: content-defined
chunking (Rabin 1981 / LBFS 2001) for insertion/deletion robustness.
T7.18: parallelizable encoding across $P$ workers with $O(N/P + \sqrt{N})$
span. T7.19: streaming encoder with $O(K)$ memory footer-layout
(independent-of-$N$ memory footprint).

\emph{Why it matters.} The engineering deliverables. Together they
unlock production deployment scenarios: heterogeneous data (T7.16),
edit-friendly archives (T7.17), HPC-scale encoding (T7.18), and
memory-constrained embedded encoders (T7.19).
\end{quote}

\textbf{Theorem 7.20 (cache-aware decoder amortization) ---
{[}operational{]}}

\begin{quote}
\emph{Statement.} With LRU cache of size $C$ achieving long-run hit
rate $\rho_C$ against the workload, amortized per-query decode latency
is $(1-\rho_C) \cdot K \cdot \mathrm{cost}(M')$. For static Zipf
popularity $p_k \propto k^{-\alpha}$ over $m$ items, the static-top
hit rate is $\rho_C \approx (C/m)^{1-\alpha}$ ($\alpha < 1$ divergent
normalizer) or $\rho_C \to 1 - C^{1-\alpha}/((\alpha-1)\zeta(\alpha))$
($\alpha > 1$ convergent normalizer).

\emph{Why it matters.} Quantifies the cache amortization gain
applicable to LRU caches in production decoders. Codex (round 3)
verified the Zipf asymptotic and corrected an earlier
``light/heavy-tailed'' mischaracterization (Zipf is heavy-tailed
in both regimes; the split is about normalizer convergence).
\end{quote}

\textbf{Theorem 7.21 (predictor-archive size trade-off) ---
{[}optimization{]}}

\begin{quote}
\emph{Statement.} Under Hoffmann 2022 scaling
$H_{M_L}(X) - h(X) \sim L^{-\alpha}$ with $\alpha \approx 0.34$, the
joint cost $L_{\mathrm{model}} + N \cdot H_{M_L}(X)$ minimized over
predictor size $L$ gives
$L^\ast(N) = \Theta(N^{1/(1+\alpha)})$
--- sub-linear in archive size.

\emph{Why it matters.} For an archive of $N$ bytes, the optimal
predictor size grows as $N^{0.75}$ — much slower than the archive
itself. Tells implementers how to scale models with data.
\end{quote}

\textbf{Theorem 7.22 (fail-safe store-mode worst-case) ---
{[}safety guarantee{]}}

\begin{quote}
\emph{Statement.} With per-sub-block mode flag (RNR vs store) and
$\eta$-floor AC bounded by $\log_2 \sigma$ per byte, the worst-case
total length is $L^{\mathrm{rep}}_{\mathrm{RNR}} \leq N \log_2 |\Sigma|
+ m$ (mode-flag overhead) under ANY source and ANY predictor.

\emph{Why it matters.} Defends against adversarial predictors and
incompressible sources. The user-originated insight ``RNR should
never inflate above raw'' formalized as a worst-case guarantee.
\end{quote}

\textbf{Theorem 7.23 (per-sub-block authenticated MAC) ---
{[}cryptographic integrity{]}}

\begin{quote}
\emph{Statement.} Per-sub-block MAC tag of $\tau$ bits (e.g., HMAC-SHA256
$\tau=256$ or Poly1305 $\tau=128$) adds overhead $m \cdot \tau$ bits
($\sim 0.1\%$--3\% of archive at $K \in [1024, 16384]$) with
cryptographic forgery resistance under standard PRF assumptions
(Bellare--Namprempre 2000 ASIACRYPT LNCS 1976).

\emph{Why it matters.} Production-grade integrity for compressed
archives. Verified by detecting all 100/100 single-bit tamperings
in \texttt{scripts/verify/thm\_7\_23\_integrity\_mac.py}.
\end{quote}

\textbf{Theorem 7.24 (revised: IRM Hoeffding tail bound) ---
{[}operational SLO{]}}

\begin{quote}
\emph{Statement.} Under (a) IRM workload (queries iid from popularity
distribution $\mu$), (b) static Perfect-LFU cache, (c) bounded per-miss
cost $T_{\max}$, (d) workload $\perp$ source: total batched-decode
latency satisfies plain iid Hoeffding tail
$\Pr(|T_{\mathrm{total}} - Q \cdot \bar T| > \sqrt{(Q/2)\log(2/\delta)} \cdot T_{\max}) \leq \delta$.

\emph{History.} An earlier version (LRU + iid Bernoulli hits) was
demolished by codex: workload $A,B,A,B,\ldots$ against $C=1$ LRU
produces $\Theta(Q)$ miss-count deviation, not $\Theta(\sqrt{Q})$.
Retracted (commit 7204189) and rewritten under restricted IRM
hypotheses (commit ebb8e21) that survive disproof (``Attack failed'').

\emph{Why it matters.} The cleanest demonstration in the validation
record of the 3-step protocol catching its own errors. Provides a
minimal honest production SLO tail bound; the LRU+IRM finite-batch
concentration it sidestepped is now closed by Theorem 7.26.
\end{quote}

\textbf{Theorem 7.26 (LRU+IRM miss-count concentration) ---
{[}operational SLO, general cache{]}}

\begin{quote}
\emph{Statement.} Under IRM workload with full-support popularity and
an \emph{LRU} cache, the miss count over $Q$ queries concentrates with
$\sqrt{Q}$-Hoeffding scaling, via the Lezaud--Gillman Markov-chain
Chernoff bound on the augmented cache-state chain
$(\mathrm{Cache}_t, X_t)$ (strictly positive spectral gap
$\gamma_{\rm LRU} > 0$, with $\gamma_{\rm LRU} \geq \pi_{\min}$ from
Dobrushin coupling).

\emph{Proof.} Lezaud 1998 / Gillman 1998 Chernoff bound for
finite-state Markov chains applied to the augmented cache-state chain.

\emph{Why it matters.} Generalises T7.24 from static Perfect-LFU to
the deployed LRU policy, closing the open question from earlier drafts
of §13.2.
\end{quote}

\textbf{Theorem 7.27 (space-time product converse) --- {[}lower
bound, tight{]}}

\begin{quote}
\emph{Statement.} For the sub-block random-access class (Def 10.3),
per-query work $C(K)$ and seek-index storage $S(K)$ obey
$C(K)\cdot S(K) \geq (1-o(1))\,N\,\mathrm{cost}(M)\log_2 N$, matching
T7.5's achievability at $K=\Theta(\sqrt N)$. The general-scheme
analogue is \emph{false} (T7.27$'$: a Ferragina--Venturini
$O(1)$-access escape attains Shannon rate for finite-order-Markov
sources), but Corollary 7.27c recovers the product on the
predictor-faithful class (covers state-snapshot schemes).

\emph{Proof.} Sync-index storage lower bound (Lemma 7.27b) $\times$
sequential per-query predictor-evaluation cost; AM--GM balance.

\emph{Why it matters.} The converse companion to T7.5: within the
sub-block class the $\sqrt N$ trade-off is not just optimal-by-design
but information-theoretically forced. The honest scope (sub-block,
not general) is established by the T7.27$'$ refutation.
\end{quote}

\textbf{Theorem 8.1 (no universal dominance) --- {[}impossibility{]}}

\begin{quote}
\emph{Statement.} No injective map from $\{0,1\}^n$ to binary
strings of length less than $n$ assigns every length-$n$ input a
strictly shorter codeword.

\emph{Proof.} Pigeonhole: $2^n$ inputs, only $2^n - 1$ shorter
binary strings exist.

\emph{Why it matters.} Standard. Included to fix the honest claim: RNR
is conditionally, not universally, advantageous.
\end{quote}

\textbf{Theorem 9.1 (no cross-subblock error propagation) --- {[}safety
guarantee{]}}

\begin{quote}
\emph{Statement.} If verification $H_v$ is a cryptographic hash, the
probability that an incorrect decoded subblock is used to update the
next-block context is at most the collision probability $\delta$.

\emph{Proof.} Context-update rule is conditional on hash match;
mismatch with $\delta$-collision-resistant hash is $\delta$-rare.

\emph{Why it matters.} Defends autoregressive correctness. With a
64-bit truncated SHA-256, $\delta \approx 2^{-64}$; cross-block
error propagation is operationally impossible.
\end{quote}

\textbf{Theorem 10.1 (bit-exact forward pass) --- {[}reproducibility{]}}

\begin{quote}
\emph{Statement.} Under (i)~integer arithmetic with sufficient
bit-width, (ii)~deterministic non-linearity tables, (iii)~fixed
reduction order, $M(C)$ is bit-identical across any pair of
conforming implementations.

\emph{Proof.} Integer arithmetic has no platform variance; deterministic
tables and fixed order eliminate associativity and selection issues;
induction on layer depth completes the proof.

\emph{Why it matters.} Reduces the reproducibility question to a
conformance specification. Required for any deployable lossless RNR
codec.
\end{quote}

\textbf{Definition 10.1 (compatible online optimizer) ---
{[}definition{]}}

\begin{quote}
\emph{Statement.} An online optimizer is compatible if its state is
integer-valued, its update is a deterministic integer-to-integer
function, and any non-integer operations are replaced by deterministic
integer approximations.

\emph{Examples.} Integer SGD, signSGD, integer Adam with reciprocal-
square-root replaced by an integer-table lookup.

\emph{Why it matters.} Names the class of optimizers that admit
bit-exact reproduction. Required input to Theorem 10.2.
\end{quote}

\textbf{Theorem 10.2 (bit-exact adaptive predictors) ---
{[}reproducibility, adaptive{]}}

\begin{quote}
\emph{Statement.} If $M$ satisfies Theorem 10.1 and adaptation uses
a compatible online optimizer with deterministic integer
backpropagation and pinned subblock order, then the encoder and
decoder's parameter states are bit-identical at every adaptation
step.

\emph{Proof.} Induction on adaptation step: bit-identical state plus
bit-identical gradient plus deterministic integer-to-integer update
yields bit-identical next state. Verified subblocks (Theorem 9.1) ensure
inputs match.

\emph{Why it matters.} Extends 10.1 to online-adaptive predictors. With
published low-precision-training methods, online adaptation is
theoretically reproducible; the engineering work is to pick a specific
compatible optimizer.
\end{quote}

\textbf{Definition 10.2 ($W$-bounded predictor) --- {[}definition{]}}

\begin{quote}
\emph{Statement.} $M$ is $W$-bounded if $M(C_i)$ depends only on the
previous $W$ bytes of public context.

\emph{Examples.} Transformers with fixed context window, $n$-gram and
finite-state models, any feedforward predictor with bounded receptive
field. Not satisfied by RNNs, LSTMs, or state-space models.

\emph{Why it matters.} Enables the random-access theorem. The
bounded-context property is what makes seek-and-read work with cost
polylogarithmic in archive size.
\end{quote}

\textbf{Definition 10.3 (Sync point) --- {[}definition{]}}

\begin{quote}
\emph{Statement.} A position $P$ where the predictor's context is
reset to a known initial value, with the encoder maintaining a seek
index from sync point identifier to bit offset in $R$.

\emph{Construction.} Encoder inserts sync points every $K$
positions and writes a per-sync index in the archive footer.

\emph{Why it matters.} The seek-and-read mechanism. Sync-point
spacing $K$ trades off random-access speed against the cold-context
compression cost.
\end{quote}

\textbf{Theorem 10.3 (random access in $O(K + k)$ time) ---
{[}seek-and-read{]}}

\begin{quote}
\emph{Statement.} For $W$-bounded $M$ with sync points every $K$
positions, random access to $k$ bytes at any position $N$ requires
$O(\log(N/K) + (K + k) \cdot \mathrm{cost}(M))$ time, only
polylogarithmic in $N$.

\emph{Proof.} Binary-search the seek index for the nearest sync
point $P$ in $O(\log(N/K))$; decode forward from $P$ to $N + k$ in
$O((K + k) \cdot \mathrm{cost}(M))$; output the requested bytes.

\emph{Why it matters.} The kill feature relative to other neural
compressors. NNCP, PixelCNN++, and Bit-Swap require $O(N)$ time per
access; RNR achieves $O(\log(N/K) + (K + k) \cdot \mathrm{cost}(M))$
at approximately 1\% ratio overhead.
Matches bgzip and zstd seekable for speed while preserving neural
ratio.
\end{quote}

\textbf{Theorem 10.4 (random access for non-$W$-bounded predictors)
--- {[}seek-and-read, general{]}}

\begin{quote}
\emph{Statement.} For non-$W$-bounded $M$ (RNN, LSTM, state-space),
random access requires storing $M$'s full internal state at each
sync point. Cost remains $O((K + k) \cdot \mathrm{cost}(M))$ per
access; seek index grows by $L(\text{state})$ bits per sync point.

\emph{Proof.} Augment each sync point with a serialized state
snapshot; deserialize on seek; decode forward as in Theorem 10.3.

\emph{Why it matters.} Quantifies the deployment penalty of using
recurrent predictors. For LSTM with $d = 1024$, snapshots add
roughly 3\% to the seek index; for state-space models, can double
the index size. Motivates the $W$-bounded design choice for
deployable codecs.
\end{quote}

\textbf{Theorem 10.5 (block-device interface) --- {[}deployment{]}}

\begin{quote}
\emph{Statement.} An OS-level block-device driver mapping
LBA-indexed reads to RNR random-access calls serves any read request
in $O(\log(N/K) + (K + k) \cdot \mathrm{cost}(M))$ time and
$O(K + k)$ working memory, with only polylogarithmic dependence on
archive size $N$.

\emph{Proof.} Direct corollary of Theorem 10.3 in driver terms:
translate $(\mathrm{LBA}, k)$ to $(\text{byte\_position}, k)$ and
call the range-decode function. No per-archive scaling.

\emph{Why it matters.} Enables the squashfs / EROFS /
btrfs-with-zstd deployment pattern with a neural compressor in the
base layer. Existing neural compressors cannot serve this use case
at all.
\end{quote}

\textbf{Theorem 10.6 (hardware-target portability) ---
{[}deployment, portability{]}}

\begin{quote}
\emph{Statement.} Any hardware target $T$ providing bit-exact
integer arithmetic at the specified bit-widths and lookup-table
evaluation produces bit-identical output to any other conforming
target. Choice of $T$ affects throughput and latency but not the
produced bytes.

\emph{Proof.} Direct corollary of Theorem 10.1 applied to specific
hardware classes. Integer arithmetic is associative and commutative
on any platform implementing two's-complement integers, which is
universal.

\emph{Why it matters.} Conforming targets include CPU scalar on any
ISA, CPU SIMD (AVX2, AVX-512, NEON, RVV), GPU integer Tensor Cores
(NVIDIA Turing+, AMD CDNA), integer NPUs (Apple ANE, Google Edge
TPU, Qualcomm Hexagon, AWS Inferentia). Encode on server GPU, decode
on phone NPU, byte-identity guaranteed.
\end{quote}

\textbf{Quick triage: reviewer challenges → answer}

\emph{For each common reviewer challenge, the relevant theorem or remark
to cite. Memorize this table for defense.}

\begin{quote}
\textbf{Q:} \emph{``Why do you think Type-I is competitive at all?''}

\textbf{A:} Theorem 4.1 (achievability), Theorem 4.2 (matching
converse).

\textbf{Q:} \emph{``Could a better algorithm find a better Code12?''}

\textbf{A:} Theorems 4.3 and 4.4: optimal selection is NP-hard and
inapproximable.

\textbf{Q:} \emph{``Your earlier draft inflated the Type-II gain.''}

\textbf{A:} Theorem 5.3 (corrected) uses multinomial baseline and
explicit escape cost.

\textbf{Q:} \emph{``Does Type-II actually do anything Type-I cannot?''}

\textbf{A:} Theorem 5.4 (linear-code separation, $\Theta(N)$).

\textbf{Q:} \emph{``For what data classes is your separation
actually realized?''}

\textbf{A:} Theorem 5.5 ($\mathrm{TC}$ is the exact quantity);
Remark 5.6 (manifold hypothesis predicts
$\mathrm{TC} = \Theta(N)$ for natural data).

\textbf{Q:} \emph{``Doesn't arithmetic coding under $Q$ dominate
your support-set scheme?''}

\textbf{A:} Arithmetic coding under $Q$ achieves cross-entropy
$\mathbb{E}_D[-\sum_i \log_2 Q_{X_i}(i)]$, which equals $H(X \mid Q)$
only when $Q$ matches the true conditional law. The optimal
fixed-size support is currently an open problem (Remark 5.7); the
natural top-$Q$-prefix heuristic is not provably optimal. Support-set
remains preferred for hard-candidate-set predictors and for hardware
regimes where rank coding is cheaper than arithmetic streaming.

\textbf{Q:} \emph{``Why $2\varepsilon$ and not $\varepsilon$?''}

\textbf{A:} Theorem 6.2: $2\varepsilon$ is tight via a constructive
adversarial example.

\textbf{Q:} \emph{``Why use byte-tensor instead of separate nibble
coders?''}

\textbf{A:} Theorem 6.3: the gain equals $I(H ; L \mid C)$, which
is sizable for ASCII text, opcode streams, and pixel data.

\textbf{Q:} \emph{``Is the choice of decomposition tree material?''}

\textbf{A:} Theorem 6.4: no, all rooted binary tree decompositions
achieve $H(B \mid C)$; it's a calibration question, not a rate
question.

\textbf{Q:} \emph{``Why can't BPE just be optimized further?''}

\textbf{A:} Theorems 6.7 and 6.8: dictionary learning is NP-hard
and APX-hard via SGP.

\textbf{Q:} \emph{``Bits-back is not really a special case of your
framework.''}

\textbf{A:} Theorem 6.6: it is, with exact KL accounting in the
rate decomposition.

\textbf{Q:} \emph{``Your framework can't be deployed because
floating-point isn't reproducible.''}

\textbf{A:} Theorems 10.1 and 10.2: a precision-and-determinism
specification suffices for both fixed and adaptive predictors;
published optimizers fit Definition 10.1.

\textbf{Q:} \emph{``Can I seek into a compressed archive without
decoding everything?''}

\textbf{A:} Theorem 10.3: yes, in $O(K + k)$ time independent of
archive size, for $W$-bounded predictors (Definition 10.2) with
sync points (Definition 10.3). Compression cost approximately 1
percent. NNCP and PixelCNN-class systems cannot do this. Theorem
10.4 covers RNN/LSTM/state-space variants via state snapshots.

\textbf{Q:} \emph{``Can I mount a compressed disk image as a filesystem
and use it like a real drive?''}

\textbf{A:} Theorem 10.5: yes. An OS-level driver serves block-device
read requests with cost polylogarithmic in archive size. Use the
squashfs+overlayfs pattern: read-only RNR base, copy-on-write writable
layer. The driver architecture is identical to existing compressed
filesystems except the predictor replaces the dictionary decompressor.

\textbf{Q:} \emph{``Can I run the decoder on a GPU / NPU / AVX-only CPU
and get the same bytes?''}

\textbf{A:} Theorem 10.6: yes. Bit-exact integer arithmetic and lookup
tables conform on any target supporting them, which is universal across
CPU scalar, CPU SIMD (AVX2, AVX-512, NEON), GPU integer Tensor Cores,
and integer NPUs (Apple ANE, Edge TPU, Hexagon, Inferentia). Choice of
target affects latency but not output.

\textbf{Q:} \emph{``What about cross-subblock errors corrupting the
context?''}

\textbf{A:} Theorem 9.1: cryptographic-hash verification gates
context updates; propagation is $\delta$-rare.

\textbf{Q:} \emph{``You're claiming universal dominance over LZMA.''}

\textbf{A:} Theorem 8.1: we explicitly disclaim universal dominance and
only assert conditional advantage.
\end{quote}

\textbf{Honest qualifications to keep ready}

\begin{itemize}
\item Theorem 4.4's reduction relies on QAP inapproximability and
an isometric Hamming-embedding argument; a fully self-contained
proof would expand the embedding step.

\item Theorem 5.5 is unconditional; Remark 5.6's quantification
depends on the manifold hypothesis being operationalized for the
specific data class. Empirical intrinsic-dimension estimation is out
of scope (§13.3).

\item Theorem 6.2 is tight in worst case; typical-case behavior is
much better. The $2\varepsilon$ factor is conservative.

\item Theorem 6.8 inherits an APX-hardness constant from SGP; the
exact constant is small and may not be practically meaningful, only
complexity-theoretically.

\item Theorem 10.1 assumes the model description specifies the
reduction order. This is a constraint on the spec, not a deep
mathematical assumption.

\item Theorem 10.2 covers any compatible optimizer; whether a
specific compatible Adam variant converges as well as
floating-point Adam on standard benchmarks is an empirical question
(open in §13.2).

\item The framework is a math paper. It does not contain empirical
results. Numbers in §11 of the main paper are explicitly labeled
as conjectures.
\end{itemize}

\end{document}
